% GENERATED FILE. Edit canonical lessons, then run npm run generate:books.
% canonical-source-hash: fnv1a64:b326926eabefc8d6
% canonical-lessons: UR-C235-kira, UR-C235-jugnu, UR-C235-khatmal, UR-C235-kekra, UR-C235-jhinga

\chapter{Insect, Firefly, Bedbug, Crab, Prawn}
\label{ch:ur-insect-firefly-bedbug-crab-prawn}
\hlchaptermodality{235}

\begin{chapteropening}
\textbf{By the end of this chapter:} \emph{I can say an insect, a firefly, a bedbug, a crab and a prawn.}

Complete the last lesson of chapter 235: I can say an insect, a firefly, a bedbug, a crab and a prawn.
\end{chapteropening}

\section[\texorpdfstring{\ur{کیڑا} (kīṛā)}{کیڑا (kīṛā)}]{\ur{کیڑا} (kīṛā) — an insect, a worm}

\label{lesson:UR-C235-kira}



\begin{quote}
Before the new one: say the Urdu for nihari, then the Urdu for minced meat.
\end{quote}

\subsection*{You'll want to know: \ur{کیڑا}}
\textbf{\ur{کیڑا}} — \emph{kīṛā} — \textquotedblleft{}an insect, a worm\textquotedblright{}.

\begin{cousinweb}
\textbf{\ur{کیڑا}} (\emph{kīṛā}) is a tadbhava, a word inherited from Sanskrit and reshaped by sound change. It comes from \emph{kīṭa}, \textquotedblleft{}insect, worm\textquotedblright{} (CDIAL 3193).
\end{cousinweb}

\begin{grammarlens}[title={What you've built}]
One more word for this chapter.
\end{grammarlens}

\subsection*{Your turn}
\begin{itemize}
\raggedright
  \item \emph{Say it:} \emph{kīṛā}
  \item \emph{Say it:} \emph{kīṛā}, once more
  \item \emph{Say it:} say \emph{kīṛā}
  \item \emph{Recall:} say the Urdu for nihari, then the Urdu for minced meat, then say \emph{kīṛā} again
\end{itemize}

\begin{tcolorbox}[breakable,colback=teal!4,colframe=teal!35!black,title={Before you move on}]
What does \ur{کیڑا} mean? (\textquotedblleft{}An insect, a worm\textquotedblright{}.) Say it once more.
\end{tcolorbox}
\section[\texorpdfstring{\ur{جگنو} (jugnū)}{جگنو (jugnū)}]{\ur{جگنو} (jugnū) — a firefly}

\label{lesson:UR-C235-jugnu}



\begin{quote}
Before the new one: say the Urdu for minced meat, then the Urdu for an insect.
\end{quote}

\subsection*{You'll want to know: \ur{جگنو}}
\textbf{\ur{جگنو}} — \emph{jugnū} — \textquotedblleft{}a firefly\textquotedblright{}.

\begin{cousinweb}
\textbf{\ur{جگنو}} (\emph{jugnū}) is a tadbhava, a word inherited from Sanskrit and reshaped by sound change. It comes from \emph{jyotiriṅgaṇa}, \textquotedblleft{}firefly\textquotedblright{} (CDIAL 5298). It is built from \emph{jyotis}, \textquotedblleft{}light\textquotedblright{}, and \emph{iṅgaṇa}, \textquotedblleft{}moving\textquotedblright{}.
\end{cousinweb}

\begin{grammarlens}[title={What you've built}]
One more word for this chapter.
\end{grammarlens}

\subsection*{Your turn}
\begin{itemize}
\raggedright
  \item \emph{Say it:} \emph{jugnū}
  \item \emph{Say it:} \emph{jugnū}, once more
  \item \emph{Say it:} say \emph{jugnū}
  \item \emph{Recall:} say the Urdu for minced meat, then the Urdu for an insect, then say \emph{jugnū} again
\end{itemize}

\begin{tcolorbox}[breakable,colback=teal!4,colframe=teal!35!black,title={Before you move on}]
What does \ur{جگنو} mean? (\textquotedblleft{}A firefly\textquotedblright{}.) Say it once more.
\end{tcolorbox}
\section[\texorpdfstring{\ur{کھٹمل} (khaṭmal)}{کھٹمل (khaṭmal)}]{\ur{کھٹمل} (khaṭmal) — a bedbug}

\label{lesson:UR-C235-khatmal}



\begin{quote}
Before the new one: say the Urdu for an insect, then the Urdu for a firefly.
\end{quote}

\subsection*{You'll want to know: \ur{کھٹمل}}
\textbf{\ur{کھٹمل}} — \emph{khaṭmal} — \textquotedblleft{}a bedbug\textquotedblright{}.

\begin{cousinweb}
\textbf{\ur{کھٹمل}} (\emph{khaṭmal}) is usually explained as \emph{khāṭ}, \textquotedblleft{}bedstead\textquotedblright{}, with \emph{mal}, \textquotedblleft{}dirt\textquotedblright{}: the pest of the bed. CDIAL lists \emph{khāṭ} under Sanskrit \emph{khaṭvā}, \textquotedblleft{}bedstead\textquotedblright{} (CDIAL 3781), and \emph{mal} under \emph{mala}, \textquotedblleft{}dirt\textquotedblright{} (CDIAL 9899), but it has no entry for \emph{khaṭmal} itself, so the account is not certain.
\end{cousinweb}

\begin{grammarlens}[title={What you've built}]
One more word for this chapter.
\end{grammarlens}

\subsection*{Your turn}
\begin{itemize}
\raggedright
  \item \emph{Say it:} \emph{khaṭmal}
  \item \emph{Say it:} \emph{khaṭmal}, once more
  \item \emph{Say it:} say \emph{khaṭmal}
  \item \emph{Recall:} say the Urdu for an insect, then the Urdu for a firefly, then say \emph{khaṭmal} again
\end{itemize}

\begin{tcolorbox}[breakable,colback=teal!4,colframe=teal!35!black,title={Before you move on}]
What does \ur{کھٹمل} mean? (\textquotedblleft{}A bedbug\textquotedblright{}.) Say it once more.
\end{tcolorbox}
\section[\texorpdfstring{\ur{کیکڑا} (kekṛā)}{کیکڑا (kekṛā)}]{\ur{کیکڑا} (kekṛā) — a crab}

\label{lesson:UR-C235-kekra}



\begin{quote}
Before the new one: say the Urdu for a firefly, then the Urdu for a bedbug.
\end{quote}

\subsection*{You'll want to know: \ur{کیکڑا}}
\textbf{\ur{کیکڑا}} — \emph{kekṛā} — \textquotedblleft{}a crab\textquotedblright{}.

\begin{cousinweb}
\textbf{\ur{کیکڑا}} (\emph{kekṛā}) is a tadbhava, a word inherited from Sanskrit and reshaped by sound change. It comes from \emph{karkaṭa}, \textquotedblleft{}crab\textquotedblright{} (CDIAL 2816). CDIAL thinks the word probably non-Aryan, and doubts the usual link with Latin \emph{cancer}.
\end{cousinweb}

\begin{grammarlens}[title={What you've built}]
One more word for this chapter.
\end{grammarlens}

\subsection*{Your turn}
\begin{itemize}
\raggedright
  \item \emph{Say it:} \emph{kekṛā}
  \item \emph{Say it:} \emph{kekṛā}, once more
  \item \emph{Say it:} say \emph{kekṛā}
  \item \emph{Recall:} say the Urdu for a firefly, then the Urdu for a bedbug, then say \emph{kekṛā} again
\end{itemize}

\begin{tcolorbox}[breakable,colback=teal!4,colframe=teal!35!black,title={Before you move on}]
What does \ur{کیکڑا} mean? (\textquotedblleft{}A crab\textquotedblright{}.) Say it once more.
\end{tcolorbox}
\section[\texorpdfstring{\ur{جھینگا} (jhīṅgā)}{جھینگا (jhīṅgā)}]{\ur{جھینگا} (jhīṅgā) — a prawn}

\label{lesson:UR-C235-jhinga}



\begin{quote}
Before the new one: say the Urdu for a bedbug, then the Urdu for a crab.
\end{quote}

\subsection*{You'll want to know: \ur{جھینگا}}
\textbf{\ur{جھینگا}} — \emph{jhīṅgā} — \textquotedblleft{}a prawn\textquotedblright{}.

\begin{cousinweb}
\textbf{\ur{جھینگا}} (\emph{jhīṅgā}) is a tadbhava, a word inherited and reshaped by sound change. It comes from \emph{*jhiṅga}, \textquotedblleft{}prawn\textquotedblright{}, a form CDIAL reconstructs from its descendants (CDIAL 5385).
\end{cousinweb}

\begin{grammarlens}[title={What you've built}]
One more word for this chapter.
\end{grammarlens}

\subsection*{Your turn}
\begin{itemize}
\raggedright
  \item \emph{Say it:} \emph{jhīṅgā}
  \item \emph{Say it:} \emph{jhīṅgā}, once more
  \item \emph{Say it:} say \emph{jhīṅgā}
  \item \emph{Recall:} say the Urdu for a bedbug, then the Urdu for a crab, then say \emph{jhīṅgā} again
\end{itemize}

\begin{tcolorbox}[breakable,colback=teal!4,colframe=teal!35!black,title={Before you move on}]
What does \ur{جھینگا} mean? (\textquotedblleft{}A prawn\textquotedblright{}.) Say it once more.
\end{tcolorbox}
